
En
notant @@M0@@ le sous-module de @@M1@@ engendré par @@M2@@, on voit donc que l'image de @@M3@@ est @@M4@@. Par ailleurs, @@M5@@ est régulier dans @@M6@@, car sur @@M7@@ l'idéal inversible @@M8@@ est défini par l'équation @@M9@@. Le noyau de @@M10@@ est donc @@M11@@. Par suite, le @@M12@@-ième groupe d'homologie de @@M13@@ est @@M14@@; on trouve donc sur @@M15@@ des isomorphismes:
@@M16@@
En particulier, on trouve pour @@M17@@:
@@M18@@
Par suite, on a pour tout @@M19@@:
@@M20@@

Or d'après 2.2, l'isomorphisme:
@@M21@@
défini par les isomorphismes (3.2.2) est un isomorphisme d'algèbres. Il en est donc de même pour l'isomorphisme:
@@M22@@
défini par les isomorphismes (3.2.3). Pour des raisons d'unicité, c'est alors la restriction à @@M23@@ de l'homomorphisme (3.2.1), et ce dernier est donc un isomorphisme.

ii) On considère la suite exacte:
@@M24@@
